\chapter{Typography}
\label{cha:typography}
\epigraph{Typography is the craft of endowing human language with a durable
visual form.}{Robert Bringhurst}

The body is Inter at 10\,pt on a 1.20 leading, which puts about 70 characters on a line. Headings are Lexend, which descends in weight rather than only in size — a section is set at 600, a
subsection at 500, and a sub-subsection drops to Inter entirely. That
progression is what makes a page feel like a book instead of a slide deck.

\section{Weights and emphasis}

Both families are variable fonts, so the weights here are real instances
pulled off the \texttt{wght} axis, not synthesised. \textbf{Bold} sits at 650,
which holds up better in print than a true 700. \emph{Italic} comes from
Inter's own italic cut.

The display cuts are available directly when you need them:
{\displaylight Lexend Light}, {\displaybook Book}, {\displaymed Medium},
{\displaysemi SemiBold}, {\displaybold Bold}.

\subsection{Inline furniture}

Press \kbd{Ctrl}+\kbd{Shift}+\kbd{P} to open the palette, then run
\cmd{git rebase --interactive}. This release is \pill[green]{stable}; the
old API is \pill[red]{deprecated} and the streaming path is
\outlinepill[teal]{beta}. Things worth \hl{marking as you read} stay marked.

Tags work the way they do in Obsidian: \tagpill{\#linear-algebra}
\tagpill[purple]{\#unfinished} \tagpill[teal]{\#reference}.

\subsubsection{A quieter heading}

Sub-subsections are set in Inter SemiBold at reading size, in the muted ink.
They are meant to organise a long passage without interrupting it.

\paragraph{Run-in headings} are useful for definition-dense notes, where a
full heading would cost more vertical space than the idea is worth.

\section{The margin}

The outer margin is 34\,mm, which is wide on purpose.\sidenote{Margin notes live
here. Keep them to a sentence or two — anything longer belongs in the body or
a callout.} It holds margin notes and gives the text block somewhere to
breathe.

\divider

Footnotes are small and quiet, with a short rule above them.\footnote{Like
this one. Set at 7.4\,pt, hanging, pinned to the bottom of the page.}

\begin{quotation}
A pull-quote is unboxed and italic. It borrows emphasis from the surrounding
whitespace rather than from a frame.
\end{quotation}

\section{Mathematics}

Maths is set in Latin Modern Math, which pairs acceptably with Inter and can
be swapped in \texttt{styles/fonts.sty} for STIX Two Math or Libertinus Math.

\begin{definition}[Convergence]
A sequence $(x_n)$ in a metric space $(X, d)$ converges to $x \in X$ if for
every $\varepsilon > 0$ there is an $N$ such that $d(x_n, x) < \varepsilon$
for all $n > N$.
\end{definition}

\begin{theorem}[Uniqueness of limits]
A sequence in a Hausdorff space has at most one limit.
\end{theorem}

\begin{proof}
Suppose $x_n \to x$ and $x_n \to y$ with $x \neq y$. Pick disjoint
neighbourhoods $U \ni x$ and $V \ni y$. The tail of the sequence lies in both,
which is impossible since $U \cap V = \emptyset$.
\end{proof}

Display maths gets the space it needs:
\[
  \int_{-\infty}^{\infty} e^{-x^2}\,\mathrm{d}x = \sqrt{\pi}.
\]
